% ---
% id: meat-core/frank-wolfe/lem:fo-face
% kind: lem
% title: "Every direction is admissible at maximal support"
% summary: >-
%   The image of every Hermitian input direction under a branch map lies in
%   the maximal branch support, so at a branch operator with maximal support no
%   direction leaves its support. The face derivative then covers every
%   feasible direction, and no one-sided derivative is infinite along them.
% prominence: standard
% status: implemented
% uses: [def:first-order-setting, def:relaxed-feasible-set, meat-core/thm:choi, meat-core/cor:full-gradient, meat-core/def:branch-support, meat-core/lem:branch-support, meat-core/thm:gradient, meat-core/lem:branch-maps]
% mentions: [thm:firstorder]
% assumes: []
% refs: []
% provenance: newly-derived
% implements: certiqs_key.meat_core.objective.gradient:z_invariant_face
% lean:
% ---

Theorem~\ref{thm:firstorder} pairs the gradient of the objective against
every feasible direction. The lemma below is what makes maximal support
of the branch operators enough for the linearisation to exist along every
direction the theorem pairs against: a \tref{branch map} sends no
Hermitian direction outside the \tref{maximal branch support}.

\begin{lemma}\label{lem:fo-face}
In the setting of Definition~\ref{def:first-order-setting}, let
$\hat c\in\Chkey$ and $s\in\SecSet_K(\hat c)$, let
$J_s\in\Pos{A_0\Qsys_s}$, and put $X\coloneqq\cG^{J}_{\hat c,s}(J_s)$.
Then for \emph{every} Hermitian $\Delta$ on $\Hs{A_0\Qsys_s}$ the
operator $H\coloneqq\cG^{J}_{\hat c,s}(\Delta)$ satisfies
\begin{equation}\label{eq:fo-face}
  \operatorname{supp}H\;\subseteq\;W_{\hat c,s}
\end{equation}
with $W_{\hat c,s}$ the maximal branch support \eqref{eq:Wcs}.
Consequently, if $X$ attains maximal support,
$\operatorname{supp}X=W_{\hat c,s}$, for which $J_s\succ0$ suffices,
then $\operatorname{supp}H\subseteq\operatorname{supp}X$ for every such
$\Delta$, so that:
\begin{enumerate}[label=(\alph*),leftmargin=2.2em]
\item Theorem~\ref{thm:gradient}\,(i) applies to the pair $(X,H)$, and
  $t\mapsto T(X+tH)$ is differentiable at $t=0$, two-sidedly, the
  hypothesis being symmetric in $H\mapsto-H$;
\item the $+\infty$ case of Theorem~\ref{thm:gradient}\,(ii) is
  \emph{unreachable} from $X$ along the image of any Hermitian direction
  $\Delta$ for which $X+t H\succeq0$ for small $t>0$. That is the class
  that theorem is stated on and the class along which $T$ is defined at
  all, and it contains every direction $y-x$ that
  Theorem~\ref{thm:firstorder} pairs against, since
  $X+t\cG^{J}_{\hat c,s}(J'_s-J_s)=\cG^{J}_{\hat c,s}\big((1-t)J_s+tJ'_s\big)
  \succeq0$ for $t\in[0,1]$ whenever $J_s$ and $J'_s$ are positive
  semidefinite. No such direction strictly enlarges the support of $X$,
  because by \eqref{eq:fo-face} none leaves
  $W_{\hat c,s}=\operatorname{supp}X$.
\end{enumerate}
\end{lemma}

\begin{proof}
Write $\Delta=\Delta_{+}-\Delta_{-}$ for the Jordan decomposition, with
$\Delta_{\pm}\succeq0$ supported on orthogonal subspaces. The map
$\cG^{J}_{\hat c,s}=\cG_{\hat c,s}\circ\chi_{\Qsys_s}$ is a composition of
two completely positive maps, $\chi_{\Qsys_s}$ being the conjugation
$\Delta\mapsto(W_{\tau}\otimes\idon{})\Delta(W_{\tau}^{\dagger}\otimes
\idon{})$ of Theorem~\ref{thm:choi} and $\cG_{\hat c,s}$ being completely
positive by Lemma~\ref{lem:branch-maps}\,(i); hence it is completely
positive and in particular positive: $\cG^{J}_{\hat c,s}(\Delta_{\pm})\succeq0$.
Each of $\chi_{\Qsys_s}(\Delta_{\pm})$ lies in $\Pos{A\Qsys_s}$ and is of
the form Lemma~\ref{lem:branch-support}\,(i) considers, so
\eqref{eq:Wcs-bound} gives
$\operatorname{supp}\cG^{J}_{\hat c,s}(\Delta_{\pm})\subseteq W_{\hat c,s}$.
For positive semidefinite $A,B$ one has
$\operatorname{supp}(A-B)\subseteq\operatorname{supp}A
+\operatorname{supp}B$, since a vector orthogonal to both supports is
annihilated by $A$ and by $B$ and hence by $A-B$, and the support of a
Hermitian operator is the orthogonal complement of its kernel. Applying
this to $A=\cG^{J}_{\hat c,s}(\Delta_{+})$,
$B=\cG^{J}_{\hat c,s}(\Delta_{-})$ gives \eqref{eq:fo-face}, the
right-hand side being a subspace and therefore closed under the sum. In
untrusted-source mode $\chi_{\Qsys_s}=\idmap$, $\Delta$ is Hermitian on
$\Hs{A\Qsys_s}$, and the same two steps apply with
$\Sadm=\Hs{A}$.

For (a): at maximal support \eqref{eq:fo-face} reads
$\operatorname{supp}H\subseteq\operatorname{supp}X$, which is the
hypothesis of Theorem~\ref{thm:gradient}\,(i) with $S=\operatorname{supp}X$;
that part concludes differentiability of $t\mapsto T(X+tH)$ at $t=0$, and
its hypothesis holds for $-H$ as well, so the derivative is two-sided.
For (b): Theorem~\ref{thm:gradient}\,(ii), whose hypothesis
$X+tH\succeq0$ for small $t>0$ is the scope clause of (b), locates the
$+\infty$ case at directions that \emph{strictly enlarge}
$\operatorname{supp}X$, and by \eqref{eq:fo-face} no image of a Hermitian
direction does. That $J_s\succ0$ suffices for maximal support is the
equality clause of Lemma~\ref{lem:branch-support}\,(i).
\end{proof}
